
%%%%%%
%
% $Author: Abishek $
% $Date: 12.01.2025 $
% $Path: ML24-02-Gait-Pattern-Recognition\report\Contents\en$
% $Version: 1.0 $
%
%
% !TeX encoding = utf8
% !TeX root = Rename
% !TeX TXS-program:bibliography = txs:///bibtex
%
%%%%%%

\section{From Development to Deployment}

\subsection{Data Structure}
The processed CSV file contains features such as stride length, stride time, and stride velocity. These features are labeled to differentiate between Parkinson's and normal cases. Preprocessing steps include:
\begin{itemize}
	\item Normalization of numerical features to a common scale.
	\item Removal of outliers to ensure data consistency.
	\item Splitting data into training and testing sets.
\end{itemize}

\subsection{Tools}
The following tools are used during development and deployment:
\begin{itemize}
	\item \textbf{Python}: For preprocessing data and training the TensorFlow model.
	\item \textbf{TensorFlow}: To build and train machine learning models.
	\item \textbf{Arduino IDE}: For deploying the trained model to the Arduino Nano 33 BLE Sense board.
\end{itemize}

\subsection{File Structure for Model Exchangements}
\begin{tikzpicture}[
	node distance=2cm and 3cm,
	startstop/.style={rectangle, rounded corners, minimum width=3cm, minimum height=1cm, text centered, draw=black, fill=red!30},
	process/.style={rectangle, minimum width=3cm, minimum height=1cm, text centered, draw=black, fill=blue!30},
	io/.style={trapezium, trapezium left angle=70, trapezium right angle=110, minimum width=3cm, minimum height=1cm, text centered, draw=black, fill=red!30},
	arrow/.style={thick,->,>=stealth}
	]
	
	% Nodes
	\node (input) [io] {Preprocessed CSV Data};
	\node (train) [process, below of=input] {TensorFlow Training};
	\node (save) [process, right of=train, xshift=3.5cm] {Model Conversion (.h)};
	\node (deploy) [startstop, below of=save] {Deployment on Arduino};
	
	% Arrows
	\draw [arrow] (input) -- (train);
	\draw [arrow] (train) -- node[anchor=south] {.tf, .h5} (save);
	\draw [arrow] (save) -- (deploy);
	
\end{tikzpicture}

To maintain a consistent and organized workflow, the following file structure is used:
\begin{lstlisting}[basicstyle=\ttfamily, frame=single]
	Model Exchangements Root
	|
	\|-- raw_data/         # Raw CSV files with stride data
	\|-- preprocessed_data/ # Cleaned and normalized CSV files
	\|-- trained_models/   # Models saved during TensorFlow training (.h5, .tflite)
	\|-- deployment_files/ # Deployment-ready files (.h for integration)
\end{lstlisting}

\subsection{Saving Models}
Models are saved at various stages of the workflow:
\begin{itemize}
	\item TensorFlow saves the initial training model in the \texttt{.tf} format.
	\item The model is then converted to the \texttt{.h5} format for testing and evaluation.
	\item For deployment, the model is exported as a \texttt{.h} file, which can be included directly in the Arduino IDE.
	\item Metadata, such as model version and accuracy, is stored in a JSON file for reference.
\end{itemize}

\subsection{Loading Models}
Loading models onto the Arduino Nano 33 BLE Sense involves the following steps:
\begin{itemize}
	\item Include the \texttt{.h} file in the Arduino IDE project.
	\item Use TensorFlow Lite Micro APIs to initialize and run the model on the board.
	\item Perform validation by running test inputs and comparing the outputs against expected results.
	\item Ensure that any errors in loading (e.g., due to file corruption) are detected and handled gracefully.
\end{itemize}

